% Provenance: copied from the design workspace CambridgePQ/recursion.fast, `report/sections/11-benchmark.tex`, on 2026-10-19. Read-only copy; the source is authoritative.
% 11-benchmark -- written by W3
\section{Benchmark design and overfitting}\label{sec:benchmark}

The brief asks for a benchmark of hash-chain proofs and worries that agents will overfit to it. In
this design the score is a proved bound over every R1CS in a size class (\cref{sec:cost}), so the
benchmark does not decide scores. This section explains what the benchmark is for, why the
class-uniform cost theorem rather than the security definition stops hard-coding, which circuit families are used, how
records are gated, and which designs the incentives are expected to produce.

%-----------------------------------------------------------------------------------------------
\subsection{Role of the benchmark under a class-uniform cost theorem}\label{subsec:bench-role}

The benchmark has four jobs, none of them scoring.
\begin{enumerate}[nosep,label=(\roman*)]
\item \emph{Demonstration and reproducibility.} Measured cycles, bytes and prover time on concrete
  instances are published with every record.
\item \emph{Differential test of the trusted machine.} The organiser's trusted emulator (planned,
  H-4) runs each accepting transcript, and the harness checks measured $\le$ claimed, componentwise on the cost vector. If the
  theorem is right this cannot fail. A failure points to a bug in the machine semantics, the decoder,
  the renderer or the emulator, and it halts the competition until resolved.
\item \emph{Badge and record eligibility}, which gates records on a runnable prover
  (\cref{subsec:bench-badge}).
\item \emph{Non-vacuity witness.} The trivial baseline must pass the benchmark end to end
  (\cref{sec:baseline}).
\end{enumerate}

%-----------------------------------------------------------------------------------------------
\subsection{Hard-coding: why knowledge soundness cannot stop it, and what does}\label{subsec:bench-hardcode}

Knowledge soundness is a property of each statement separately: for every $(i,x)$, acceptance implies
that an extractor finds a witness. A verifier that dispatches on the index is knowledge-sound as long
as every branch is. Such a verifier might run a cheap specialised protocol when $\idx(i)$ is the
SHA-256 chain and a general protocol otherwise. The extractor is existential and can be specialised
along with the protocol. Under a \emph{measured} score on a public hash-chain benchmark, this verifier
wins without advancing anything. Hidden instances do not help against circuit specialisation: the
circuit family is public by necessity, and only the instances vary.

The class-uniform cost theorem removes the incentive. $\mathsf{AcceptingCostAtMost}(k,c)$ quantifies over
\emph{every} index of class shape $k$. A fast path for one circuit lowers the cost on that circuit
but not the maximum, and the maximum is the score. Instance hard-coding fails for the same reason,
and also for a structural one: the verification key ($\le2^{10}$ elements) is fixed before the
instance, so it cannot hold answers to future $x$. An indexer can still precompute per-circuit data
into $\idx(i)$. That is holography, it is legitimate, and the index oracle is priced per query.

\begin{table}[t]
\centering\small
\begin{tabularx}{\linewidth}{@{}>{\raggedright\arraybackslash}p{3.2cm}>{\raggedright\arraybackslash}X>{\raggedright\arraybackslash}X@{}}
\toprule
& Proved over the class (adopted) & Measured on instances \\
\midrule
Overfitting to circuits & cannot improve the score & profitable whenever the family is public \\
Overfitting to instances & cannot improve the score & needs hidden instances \\
Prover on the scoring path & no & yes, and its choices move the score \\
Trusted base for the score & Lean machine semantics & plus emulator, generator, harness \\
Worst case or typical & worst case over indices, inputs and challenges & median or max over a sample \\
Agent burden & a cost theorem (loop invariants, closed forms in $k$) & none \\
\bottomrule
\end{tabularx}
\caption{Class-uniform cost theorem against a measured score [D2 §5.3].}
\label{tab:bench-proved}
\end{table}

What overfitting remains is to the \emph{corners of the class box} and to the \emph{cost model}, not
to the benchmark [A1-n1]. The corner case was real: with $n_x\le2^{16}$, every verifier was scored at
$2^{16}$ public inputs, and $x$ ingestion became 27--50\% of every score. v0.2 fixes $n_x=2^{10}$
(\cref{tab:cost-caps}). The cost-model residuals are listed in \cref{subsec:cost-fails}. A related
point concerns \emph{knowledge}: for $n_x=1$, a satisfiability bitmap over $K$ in the index oracle
would let a verifier accept with one query, and a brute-force extractor would make it
\enquote{knowledge-sound}. This cannot affect the score, which is worst case over $n_x=2^{10}$, but it
shows that L1 \enquote{knowledge} is nominal (\cref{subsec:sec-extractor}) [A1-n2].

\begin{decision}{BEN-1}{Benchmarks demonstrate; the theorem scores}
The score is the class-uniform accepting-run bound. Benchmark runs are unscored demonstrations, a
differential test (measured $\le$ claimed) and the gate for records. \emph{Rationale:}
special-casing cannot lower a worst case over the class, and no prover or generator sits on the
scoring path. \emph{Answers:} brief points 6 and 9 (\cref{tab:brief-review}); circuit and instance
hard-coding (D2 vector 1). [DESIGN §1, §9; D2-2, D2-8]
\end{decision}

%-----------------------------------------------------------------------------------------------
\subsection{Circuit families over \texorpdfstring{$K$}{K}}\label{subsec:bench-families}

\paragraph{Boolean R1CS of the hash compressions.} BLAKE2s and SHA-256 compression functions are
encoded as Boolean R1CS over $\mathrm{GF}(2)\subset K$, with exact counts from a symbolic generator:
\begin{itemize}[nosep]
\item Wire 0 is the constant 1. Every other wire is a bit.
\item A 32-bit modular addition has 31 product rows, one AND per bit for the carry
  $c_{i+1}=(a_i+c_i)(b_i+c_i)+c_i$, and 32 linear sum rows. The majority function is
  $\mathrm{maj}(a,b,c)=(a\oplus c)(b\oplus c)\oplus c$.
\item An XOR output is a committed wire with a linear row ($A=\sum$, $B=1$, $C=0$).
\item Rotations and shifts are renamings and cost nothing.
\item Free input bits are forced Boolean downstream by the carry rows, or by explicit
  $w\cdot w=w$ rows (1{,}280 for BLAKE2s).
\end{itemize}

\begin{table}[t]
\centering\small
\begin{tabular}{@{}lrrrrrr@{}}
\toprule
Circuit, encoding & rows & product & linear & cols & nnz & nnz/row \\
\midrule
\textbf{BLAKE2s, explicit XOR wires} & 40{,}992 & 14{,}880 & 26{,}112 & 42{,}017 & 207{,}168 & 5.05 \\
BLAKE2s, substituted (Flock style) & 15{,}136 & 14{,}880 & 256 & 16{,}161 & 82{,}267{,}098 & 5{,}435 \\
\textbf{SHA-256, explicit XOR wires} & 59{,}704 & 22{,}696 & 37{,}008 & 62{,}521 & 294{,}864 & 4.94 \\
SHA-256, substituted & 22{,}696 & 22{,}696 & 0 & 25{,}513 & 175{,}955{,}759 & 7{,}753 \\
\bottomrule
\end{tabular}
\caption{Boolean R1CS of one compression over $K$ [D3 §5.1; probe \texttt{d3-probes/circuits.py}].
The BLAKE2s product count reproduces Flock's 14{,}720 AND wires plus 160 from unfused three-operand
additions. The substituted nonzero count reproduces leanVM's \enquote{$\approx89$\,M nonzeros} for its
$2^{14}$-column block.}
\label{tab:bench-circuits}
\end{table}

Only the explicit encoding fits the class box $\mathsf{nnz}\le2^{k+4}$, i.e.\ 16 nonzeros per row on
average. The substituted encoding has 5--8k nonzeros per row. It is a prover-side trick for batches
of \emph{identical} blocks: Flock's lincheck evaluates the dense matrix by walking the circuit, which
a general-R1CS verifier cannot do. The explicit encoding is the right shape for a benchmark of
holographic verifiers.

\paragraph{Chains and classes.} A chain of $L$ compressions binds each block's output bits to the
next block's input wires. $L$ is the largest value with $\max(\text{rows},\text{cols})\le2^k$
(\cref{tab:bench-chains}).

\begin{table}[t]
\centering\small
\begin{tabular}{@{}lccc@{}}
\toprule
Family & $k=16$ & $k=20$ & $k=24$ \\
\midrule
BLAKE2s chain & $L=1$ ($2^{15.3}$ rows) & $L=24$ ($2^{19.9}$; 5.0\,M nnz) & $L=399$ ($2^{24.0}$; 82.7\,M) \\
SHA-256 chain & $L=1$ ($2^{15.9}$ rows) & $L=16$ ($2^{19.9}$; 4.7\,M nnz) & $L=268$ ($2^{23.9}$; 79.0\,M) \\
Random sparse R1CS & \multicolumn{3}{c}{$2$, $4$ or $8$ nonzeros per $A$/$B$ row; hidden seed; satisfiable by construction} \\
Boolean random (optional) & \multicolumn{3}{c}{AND/XOR DAGs as R1CS over $\mathrm{GF}(2)\subset K$} \\
\bottomrule
\end{tabular}
\caption{Benchmark families per class [D3 §5.2]. Poseidon2 chains (prime-field only) are dropped. A
binary-friendly algebraic hash over $K$ is a candidate future family.}
\label{tab:bench-chains}
\end{table}

The hash chains have about 5 nonzeros per row, three times below the class cap of 16. The
certificate's worst case lies at the cap, so the benchmark measures typical cost and the theorem
measures the worst case. The random sparse family stresses the class-uniform theorem directly: it is
unstructured, so there is nothing to special-case, and its $A$ and $B$ rows are random with the $C$
row chosen to make the instance satisfiable.

\paragraph{Public inputs and the satisfiability rule.} For a hash chain, $x$ is the first message
block and the final digest. These are packed into 40 64-bit words, whose bits index rows decompose,
so $n_x\le2^{10}$ holds (G-D3-13). This also enforces a rule from the security design: \emph{no
benchmark index may be provably always-satisfiable}. If every $x$ of an index had a witness, an
accept-all verifier would be RBR-KS on that index, with a brute-force extractor, and the benchmark
would certify nothing. Producing such a proof for a hash chain would require proving the hash
surjective onto the digest space, which is infeasible. This argument covers the hash chains, whose
public input fixes the witness and can be false. It is not made for the random sparse family, whose
$C$ rows are chosen to satisfy one instance, or for chains with free inner message blocks.

\paragraph{Padding to class shapes.} Certificates are proved only at the three class shapes:
$m=\mathrm{cols}=2^k$, $n_x=2^{10}$, $\mathsf{nnz}\le2^{k+4}$. The harness pads every smaller R1CS
with a contract-defined \lean{pad} that adds zero rows, zero witness columns and zero public inputs.
The library theorem \lean{pad\_sat\_iff} proves that \lean{pad} preserves and reflects
satisfiability (both directions, explicit witness maps), as the soundness transfer needs: padded
columns have all-zero coefficients [DESIGN §2; A3-S3]. Fixed
shapes turn WHIR and sumcheck parameters into constants, prove errors by \lean{norm\_num}, and remove
\enquote{for all sizes below the cap} casework from every proof. Small circuits pay full-class cost,
which is already the score's view. Demonstrations run on padded instances, so measured costs may fall
below the claim; the badge checks $\le$, never equality.

\paragraph{Hidden and public.} Generators, their versions and the compression R1CS are public. Before
each run, the harness draws the instances from a CSPRNG: chain preimages, and random seeds for the
random R1CS, including a fresh index that the compiled Lean indexer then preprocesses. Instances are
hidden before the run and published afterwards as (generator version, seed), so anyone can
regenerate them. Hiding does not protect the score, which is proved. It protects the badge: a native
prover tuned to a public instance list, for example one holding precomputed proofs, would otherwise
pass. Display figures are the median and maximum over 10 instances per family and class.

%-----------------------------------------------------------------------------------------------
\subsection{Record eligibility: the badge path}\label{subsec:bench-badge}

A proved score says nothing about whether the honest prover can run. The badge path checks that,
outside the theorem and without affecting the score.
\begin{enumerate}[nosep]
\item \emph{Indexer.} The harness compiles the submission's Lean indexer in a sandbox, evaluates it on
  the hidden benchmark indices, and requires the native indexer's $\vk_i$ and $\idx(i)$ to match it
  exactly.
\item \emph{Prover.} The native prover interacts round by round with challenges that the harness
  samples. On the instances up to $2^{16}$, its transcript must equal, bit for bit, that of the
  compiled Lean \lean{Strategy} on the same challenges.
\item \emph{Cost.} The trusted emulator runs every accepting transcript and checks measured $\le$
  claimed for every class.
\item \emph{Budgets.} The prover must finish within 1\,h per $2^{20}$ instance and 8\,h per $2^{24}$
  instance on the reference machine, under a stated RAM budget (A1 proposes 48\,GiB).
\end{enumerate}
Running compiled candidate Lean is acceptable here because this path decides a badge, not a theorem.
The score path executes no candidate code. \emph{Status:} these four steps are specified; the v0.3
harness runs the badge only as a hook (\code{--badge}, its own sandboxed process after the gate), and
the checks themselves, the organiser emulator and the sampler are not yet built (H-3--H-5, H-9).

\paragraph{Why step 1 is load-bearing.} $\mathsf{Agrees}$ constrains the image only on the $\vk_i$
and $\idx(i)$ that the \emph{specification} indexer computes. Demonstration runs used to take them
from the untrusted native indexer. So the following passed every check [A1-M3]:
\begin{itemize}[nosep]
\item the image accepts immediately if $\idx[0]=\texttt{0xC0FFEE}$;
\item the specification indexer always writes 0 there, so the branch is dead in the theorem and
  $\mathsf{Agrees}$ holds;
\item the native indexer writes \texttt{0xC0FFEE}, so the native prover passes every demonstration
  with no work.
\end{itemize}
The certificate is untouched, and the badge is earned by a submission whose honest prover may be
infeasible. For example, $\rho=2^{-10}$ needs $\approx128$\,GiB of $K$-valued oracles at
$2^{24}$ (more if $E$-valued). The computability gate on the badge-path fields (no
\lean{noncomputable}; no \lean{implemented\_by}, \lean{extern}, \lean{csimp}, \lean{init} or
\lean{export} entry anywhere in the solution) guarantees that \lean{VK}, \lean{IDX} and the prover
have compiled code that is the kernel's definition, and since round~4 the gate checks the build that
comparator verified (A7-M3); it does not make them feasible, since a satisfiability table is
Lean-computable by enumeration (C-7). Feasibility is what the budgets of step 4 check. The match
requirement removes the backdoor. The transcript comparison catches a native prover that departs from
the specified one.

\begin{decision}{BEN-2}{Records require the badge path}
A record needs a certified score \emph{and} a passed badge: a compiled-Lean indexer match, a
native-prover run against harness challenges, a transcript match up to $2^{16}$, measured $\le$
claimed, and wall-clock and RAM budgets. The badge is unscored. \emph{Answers:} the record-gate bypass
[A1-M3]; noncomputable indexers and provers [A2-M4]; infeasible provers hidden behind a valid
theorem (D2 vector 10). [DESIGN §8]
\end{decision}

%-----------------------------------------------------------------------------------------------
\subsection{Expected trajectory of designs}\label{subsec:bench-trajectory}

The incentives are meant to produce interesting succinct designs. \Cref{tab:bench-trajectory} gives
the binary-field trajectory model. Per-operation costs are the measured kernels of
\cref{tab:cost-field}. Protocol-level counts follow the leanVM-style binary PCS:
\begin{itemize}[nosep]
\item level-0 oracles $K$-valued, later levels $E$-valued;
\item fold factors 6 then 4;
\item 268/133/67 queries per level at $\log_2(1/\rho)=1/2/4$ in the Johnson-bound (JB) regime with
  slack $\eta=0.011$;
\item SPARK as one opening of a $2^{k+7}$-element stack, plus $3(k+4)$ degree-3 sumcheck rounds;
\item the FS surcharge and the blow-up factor of \cref{sec:cost}.
\end{itemize}
Succinct rows are accurate to about a factor of 2. \Cref{tab:bench-trajectory} prices them with the
v0.2 flat weights, which are now a display column; \cref{tab:bench-models} re-prices them under the
v0.3 models.

\begingroup\footnotesize
\setlength{\tabcolsep}{4pt}
\renewcommand{\arraystretch}{1.08}
\begin{longtable}{@{}>{\raggedright\arraybackslash}p{0.30\textwidth}rrrrrrr@{}}
\caption{Expected trajectory over $K/E$ at $\mathtt{clmul}=4$ [D3 §6.1; probe
\texttt{d3-probes/model/trajectory\_bin.py}]. $S_k$ in cycles including FS, bytes and blow-up factor;
KiB and $V$ at $2^{24}$. For scale, a leanVM 2$\to$1 step is 0.56--0.69\,M leanISA steps
and 190--206\,KiB; ported naively it would be $\approx3\times10^7$ RV64IM+Zbc cycles.}\label{tab:bench-trajectory}\\
\toprule
Design & $S_{16}$ & $S_{20}$ & $S_{24}$ & \textbf{Score} & KiB & $V$ & vs trivial \\
\midrule
\endfirsthead
\toprule
Design & $S_{16}$ & $S_{20}$ & $S_{24}$ & \textbf{Score} & KiB & $V$ & vs trivial \\
\midrule
\endhead
trivial baseline, modelled (proof $=w$), \texttt{baseline\_k.S} ($\mathsf{nnz}=16N$) & 31.6M & 506M & 8{,}095M & 506M & 3.6\,GiB & 0 & 1$\times$ \\
certified baseline (\cref{tab:base-score}, v0.2 weights) & 29.6M & 474M & 7{,}583M & 474M & -- & 0.06 & 1.07$\times$ \\
\textbf{B1} Spartan + binary WHIR/Ligerito, JB, $\rho=\tfrac12$ & 4.42M & 5.30M & 6.19M & \textbf{5.26M} & 944 & 30 & 96$\times$ \\
B2 B1 at $\rho=\tfrac14$ & 3.24M & 3.99M & 4.77M & 3.95M & 616 & 60 & 128$\times$ \\
B3 B1 at $\rho=\tfrac1{16}$ (factor 1.97) & 3.08M & 3.90M & 4.77M & 3.85M & 440 & 239 & 131$\times$ \\
B4 B1 with textbook GKR SPARK (588 rounds) & 4.99M & 6.13M & 7.34M & 6.08M & 979 & 30 & 83$\times$ \\
B5 B1, unique decoding (unconditional) & 6.74M & 8.85M & 11.21M & 8.74M & 1{,}973 & 30 & 58$\times$ \\
B7 B1 + packed Boolean witness, ring switching & 4.06M & 5.29M & 6.22M & 5.11M & 836 & 30 & 99$\times$ \\
B8 Spartan + FRI/BaseFold-shaped, $\rho=\tfrac12$ & 10.26M & 14.12M & 18.46M & 13.88M & 6{,}982 & 30 & 36$\times$ \\
B9 B8 at $\rho=\tfrac14$ & 6.56M & 8.96M & 11.63M & 8.81M & 3{,}661 & 61 & 57$\times$ \\
B10 B1 under the refuted capacity conjecture (banned) & 2.34M & 2.80M & 3.27M & 2.78M & 455 & 30 & 182$\times$ \\
\bottomrule
\end{longtable}
\endgroup

Readings:
\begin{enumerate}[nosep]
\item \emph{Succinct, proven-regime designs win by two orders of magnitude.} B1 beats the modelled
  trivial baseline 96 times and the certified one about 90 times. At the class corner the crossover is at $k\approx13$, so all three classes are
  succinct territory. No row is a cost-model exploit.
\item \emph{The levers are architectural.} B1 at $2^{24}$ spends 52\% of its 5.16\,M cycles ($c_{24}$, before the FS, byte and
  blow-up terms that bring it to $S_{24}=6.19$\,M) on the index (SPARK) opening and 41\% on the witness opening. Each query costs about 5.6k cycles: 2.4k for
  the 64-term leaf fold, 2.3k for 18 weight factors and 0.9k for 20 Merkle levels. The levers are
  therefore the proven query count (268 at $\rho=1/2$), the leaf size and initial fold, and above all
  the size of the index commitment: fewer or shorter index vectors, and sparsity exploited in the
  style of Lasso/Shout \cite{ST25-TwistShout}. That is holography engineering, the design space the
  competition exists to explore.
\item \emph{Rate is a slope.} $\rho=1/4$ gains 25\%. $\rho=1/16$ gains nothing more after the blow-up
  factor.
\item \emph{The proven-regime rule costs what it costs.} The banned capacity conjecture would gain
  1.9 times (B10), about the same as in prime fields. The unconditional (UD) track pays 1.7 times
  over JB.
\item \emph{WHIR-style folding beats FRI-shaped commitments} by 2.6 times (B8 against B1). The cause is
  15 times more queries (12{,}060 against 803 at $2^{24}$) and 7 times more proof bytes.
\item \emph{Packing is a wash.} Packed-Boolean witnesses with ring switching (B7) gain 3\%, because the
  index, not the witness, dominates.
\end{enumerate}
\paragraph{Robustness under the v0.3 cost models.} Design note D4 re-ran the trajectory with
per-model kernel costs and per-model Merkle prices (\cref{tab:bench-models}).

\begin{table}[t]
\centering\footnotesize
\setlength{\tabcolsep}{3.5pt}
\begin{tabular}{@{}lrrrrrrrr@{}}
\toprule
Design (geomean over classes, M units) & uni & zisk & sp1z & r0 & jolt & sp1 & inord & ooo \\
\midrule
trivial baseline (modelled) & 506 & 301 & 406 & 989 & 411 & 8{,}672 & 726 & 113 \\
B1 Spartan + WHIR-JB, $\rho=\tfrac12$ & 5.26 & 4.68 & 7.14 & 10.30 & 6.97 & 88.8 & 28.2 & 13.0 \\
B2 $\rho=\tfrac14$ & 3.95 & \textbf{3.58} & \textbf{5.48} & 7.66 & \textbf{5.36} & 65.0 & \textbf{22.0} & \textbf{10.2} \\
B3 $\rho=\tfrac1{16}$ & \textbf{3.86} & \textbf{3.58} & 5.49 & \textbf{7.40} & 5.40 & 61.5 & 22.5 & 10.4 \\
B4 textbook GKR SPARK & 6.09 & 5.25 & 7.86 & 11.25 & 7.58 & 101.5 & 28.7 & 13.2 \\
B5 unique decoding & 8.75 & 8.31 & 12.94 & 17.77 & 12.83 & 141.2 & 56.0 & 25.9 \\
B8 FRI/BaseFold, $\rho=\tfrac12$ & 9.37 & 15.90 & 28.96 & 27.15 & 31.62 & 96.5 & 196.8 & 91.9 \\
B9 FRI/BaseFold, $\rho=\tfrac14$ & 5.84 & 9.75 & 17.69 & 16.60 & 19.28 & \textbf{60.8} & 119.3 & 55.7 \\
\bottomrule
\end{tabular}
\caption{Trajectory under nine cost models (\texttt{zkmem} omitted; bold = best per model)
[D4 RESULT §2], with D4's pre-round-3 leaf pricing. For comparison, the \emph{certified} baseline under
the final v0.3 pricing has per-profile geomeans 844 (ZisK), 1{,}779 (SP1 with Zbc, i.e.\
\texttt{sp1z}), 2{,}358 (RISC~Zero) and 1{,}965 (Jolt), and $\mathrm{Score}_{zk}=1{,}624$
(\cref{tab:base-score}). The re-run's FRI-shaped rows differ from D3's under \texttt{uni} (9.4M and 5.8M
against 13.9M and 8.8M), within the model's $\pm2\times$ uncertainty; ranks are compared within the
re-run.}
\label{tab:bench-models}
\end{table}

\begin{itemize}[nosep]
\item \emph{The top three architectures are stable.} Spartan with binary WHIR at $\rho=1/16,1/4,1/2$
  occupies ranks 1--3 under every model except the one without Zbc. B2 and B3 swap under some models,
  but they are within 1--6\% everywhere. The optimal rate is $\log_2(1/\rho)=3$ under all nine
  models ($\rho=1/8$, not tabulated; D4).
\item \emph{FRI-shaped designs drop as soon as hashing is priced.} B9 falls from rank 4 (\texttt{uni})
  to 5--6 under every model that prices a Merkle level above 43, because it pays more levels per query.
  Without a Zbc chip (\texttt{sp1}) it is rank 1, since it has the fewest field operations. The ISA
  assumption, not the weights, decides that case.
\item \emph{The latency board ranks differently.} Under Brent's bound, $T_{16}/W$ is 0.07--0.30.
  Textbook GKR SPARK, with 588 sequential sumcheck rounds, falls from rank 4 on the work board to
  rank 7 on the latency board. FRI-shaped B9 rises from 5 to 4, and the UD regime, which has many
  parallel queries, rises from 6 to 5. The top three are unchanged. On the in-order model the two
  boards agree, because the per-query hash chain dominates both work and span.
\item \emph{Absolute levels move}: ZisK 0.9, Jolt 1.3, RISC~Zero 2.0 and the native in-order core 5.4
  times v0.2. Under native pricing, B1 is hash-dominated: about 24M of 28M.
\end{itemize}
The architectural incentive therefore survives the change of cost model. What v0.3 changes is which
\emph{engineering} pays (\cref{subsec:cost-field}) and how sequential depth is valued on the latency
board.

The expected order of improvements:
\begin{enumerate}[nosep]
\item Spartan-style zerocheck and lincheck over $K/E$ with a binary WHIR/Ligerito opening
  \cite{Set20-Spartan,ACFY24-WHIR,NA25-Ligerito}, in UD ($\approx9$\,M) or JB ($\approx5$\,M);
\item rate $1/4$ and leaf/fold tuning ($\approx4$\,M);
\item hand-written $E$ kernels ($\approx15\%$);
\item index-commitment engineering and hint-checked inversions ($\approx2$--$3$\,M);
\item accumulation (T2), with deferred index claims paid once by the decider;
\item a field-syscall board on the path to T3.
\end{enumerate}

\begin{issue}{I-BEN-1}{The leaderboard stalls without an organiser PCS}
Every succinct row needs a hash-based polynomial commitment with RBR knowledge soundness over binary
RS codes. None exists in Lean: ArkLib has no WHIR, FRI and Binius are unproved there, and leanerVM's
WHIR targets RBR soundness only. A competitive entry is 45--100k lines even with organiser libraries
[A3-M2]. Without the organiser library L4 (15--40k lines; interim tensor-code PCS 5--10k) and a
reference succinct submission L5, the leaderboard can be expected to stay at the trivial or hybrid
baseline for months. Action: staff L4 as the critical path, open with a Lean-only T0 season, and offer
the interim tensor-code PCS track (\cref{sec:issues}).
\end{issue}

%-----------------------------------------------------------------------------------------------
\subsection{Keeping the design space open}\label{subsec:bench-open}

Under one board, the leaderboard can be expected to converge within weeks to \enquote{Spartan plus
binary WHIR at the rate slope, with hand-tuned $E$ kernels} [A1-n3]. Several boards, all served by one
vector certificate, keep other architectures visible.
\begin{itemize}[nosep]
\item \emph{Unconditional and conditional proximity tracks.} UD (kernel-proven only) and JB under
  leanerVM's named \lean{McaJohnson} hypothesis (\cref{subsec:sec-hypotheses}). They have different
  optima: UD gains much less from rate than JB ($-25\%$ against $-50\%$ queries from $\rho=1/4$ to
  $1/16$, \cref{tab:bg-queries}).
\item \emph{T0 (Lean only).} The verifier is a contract query program \lean{QProg}, scored by virtual
  query cost and bytes. T0 opens first, lets protocol work start with no RISC-V proof, and carries
  into T1 through the \lean{QProg} midpoint [A3-M5, S6].
\item \emph{T1, T2 and per-class boards}, plus a Pareto display of score against $V_k$ and proof
  bytes.
\item \emph{Proof-size board.} $B_{24}$ of the decider's input, subject to $C_{24}\le4\times$ the
  current T1 record. This is the Ethereum Foundation's $\le300$\,KiB axis.
\item \emph{Per-profile columns and the latency board.} The four zkVM profile scores are shown next
  to the geometric mean, so zkVM-specific designs are visible. The latency board ranks the same
  certified submissions by $T_{16}$ (\cref{subsec:cost-lat}) and rewards parallel structure that the
  recursion board ignores. Both come from one certificate, at no extra proof cost beyond the optional
  timing clause.
\item \emph{T1-F field-syscall board.} An \texttt{EMUL} syscall priced from a binary-field zkVM
  (one leanISA step, or ZisK's native cost). There, holography, accumulation and GKR-based memory
  checking are not drowned by assembly, and it is the route towards T3 [A1-M5b].
\item \emph{Later:} a ROM track (\texttt{HASH} syscall, real Merkle and Fiat--Shamir costs,
  grinding); a Boolean track over $\mathrm{GF}(2)$ with packed witnesses; the L2 efficient-extraction
  flag; and the interim tensor-code PCS track.
\end{itemize}
